\documentclass[exercice]{thedocument}%
\usepackage{texmaths-tikz}
\startdatakeys
    \source{inconnue}
    \BO{2020}
    \cycle{4}
    \competencesDuSocle{RE}
    \connaissancesTravaillees{D2b6}
    \competencesTravaillees{A1a2, D2c4, D2a*, C2d*}
\enddatakeys
\begin{document}%
\mathspoint{A'} est l'image de \mathspoint{A} par l'homothétie $h$ de centre \mathspoint{O} et de rapport $k$.\\[0.2em]
\mathspoint{B'} est l'image de \mathspoint{B} par cette même homothétie $h$.

\begin{enumerate}
  \item Construire le centre \mathspoint{O} et déterminer le rapport $k$ de cette homothétie en justifiant.
  \item Ecrire les égalités faisant intervenir des longueurs de segment et le rapport $k$ de l'homothétie.
  \item Que peut-on dire des triangles \mathspoint{OAB} et \mathspoint{OA'B'}~?
\end{enumerate}

\begin{center}
\begin{tikzpicture}[scale=1.3]
  % Segment A'B'
  \draw[black!70, line width=0.9pt] (0,3) -- (4,5);
  % Segment AB
  \draw[black!70, line width=0.9pt] (4,0) -- (6,1);
  % Point A'
  \ptmark{0}{3}
  \node[ptcol, left, font=\small\itshape] at (0,3) {$A'$};
  % Point B'
  \ptmark{4}{5}
  \node[ptcol, above right, font=\small\itshape] at (4,5) {$B'$};
  % Point A
  \ptmark{4}{0}
  \node[ptcol, below, font=\small\itshape] at (4,0) {$A$};
  % Point B
  \ptmark{6}{1}
  \node[ptcol, right, font=\small\itshape] at (6,1) {$B$};
\end{tikzpicture}
\end{center}

\end{document}